\documentclass[conference]{IEEEtran}

\usepackage[utf8]{inputenc}
\usepackage[T1]{fontenc}
\usepackage[portuguese]{babel}
\usepackage{booktabs}
\usepackage{tabularx}
\usepackage{url}
\usepackage{hyperref}

\title{UXAgent: Auditoria de Conformidade e Relatório Técnico da Implementação Instrumentada}

\author{\IEEEauthorblockN{Projeto UXAgent}
\IEEEauthorblockA{Disciplina de Interação Humano--Computador\\
Centro Federal de Educação Tecnológica\\
Brasil}}

\begin{document}
\maketitle

\begin{abstract}
Este artigo apresenta uma auditoria da implementação do UXAgent em relação ao
enunciado da fase P3. O sistema é um agente proxy baseado em LLM que opera um
lobby cooperativo sintético, mantendo uma separação explícita entre estado
verdadeiro, observação, política, controlador, instrumentação e avaliador
externo. A implementação atende os requisitos centrais de fronteira de
informação, ações válidas por estado, mecanismos comportamentais fora do
prompt, logs JSONL validados e avaliação independente. A auditoria também
identifica pendências para a entrega completa: formalização integral de CTT e
MoLIC, matriz MAC com as categorias PP/PA/AL/RA/CE, pacote de piloto com
prompts e contextos efetivos, comparação empírica com KLM e execução do artigo
em PDF IEEE. Os resultados mock do piloto são sintéticos e não constituem
validação com participantes humanos.
\end{abstract}

\begin{IEEEkeywords}
UXAgent, agente proxy, IHC, MoLIC, CTT, KLM, instrumentação, LLM.
\end{IEEEkeywords}

\section{Introdução}
O enunciado P3 solicita a implementação de um sistema instrumentado para um
UX-Agent, incluindo projeto de interação, acoplamento com um LLM, mecanismos
parametrizados, registro estruturado, verificador externo e piloto. A pergunta
de projeto é quais decisões permitem produzir registros auditáveis fornecendo
ao agente somente a informação disponível à persona.

Este relatório responde a essa pergunta por meio de uma auditoria do
repositório no estado congelado pelo manifesto
\texttt{docs/freeze\_manifest\_v4.json}. A análise separa evidência executada,
princípios de design e hipóteses operacionais, evitando tratar resultados do
Gemini como comportamento humano.

\section{Modelagem herdada}
O sistema mantém quatro tarefas sintéticas. T1 e T2 são tarefas de calibração;
T3 e T4 são tarefas independentes reservadas para validação. A HTA refinada,
a metamensagem, os signos, os modelos CTT e GOMS/KLM, as cenas MoLIC e o
protocolo humano estão documentados em
\texttt{docs/ihc\_design\_artifacts\_v1.md} e
\texttt{docs/human\_validation\_protocol\_v1.md}.

Na T1, a política deve atribuir as funções
\texttt{LOCATE\_RESOURCE}, \texttt{ACTIVATE\_CHECKPOINT} e
\texttt{TRANSPORT\_RESOURCE}, revisar e confirmar. T2 acrescenta atualização
do feed e comunicação. T3 exige remoção de uma atribuição indisponível e
redistribuição. T4 exige integração da ordem das etapas e comunicação.

\section{Decisões de design}
\begin{table*}[t]
\caption{Decisões de design e status epistemológico}
\label{tab:decisions}
\centering
\begin{tabularx}{\textwidth}{p{0.25\textwidth}p{0.48\textwidth}p{0.17\textwidth}}
\toprule
Decisão & Justificativa & Status\\
\midrule
Separar estado verdadeiro e observação & Impede que o LLM receba o gabarito ou dados ocultos; o controlador constrói a observação e o avaliador lê o estado real. & Evidência\\
Usar vocabulário fechado e schema & Restringe a resposta a uma ação estruturada e permite rejeição auditável. & Evidência\\
Manter mecanismos R1--R4 fora do prompt & Evita encenação textual de hesitação, erro ou ajuda; os mecanismos operam no controlador. & Evidência\\
Disponibilizar somente ações autorizadas & Aproxima as falas do usuário no MoLIC e reduz ações impossíveis no estado atual. & Princípio de design\\
Registrar rejeições e rupturas por assinatura & Evita inferência a partir de discurso emocional do LLM. & Princípio de design\\
Usar T3/T4 somente como validação reservada & Evita contaminar a calibração de T1/T2. & Regra metodológica\\
Usar limites e tolerâncias preliminares & Define uma análise futura, mas ainda depende de piloto e aprovação institucional. & Hipótese operacional\\
\bottomrule
\end{tabularx}
\end{table*}

\section{Arquitetura e acoplamento}
A arquitetura possui seis camadas. \textit{Environment} mantém o estado
verdadeiro e o oráculo externo. \textit{Controller} aplica pré-condições,
transições, rejeições e mecanismos. \textit{UXAgentSession} fornece uma
fachada estável para a interface. \textit{GeminiActionPolicy} recebe uma
observação serializada e retorna uma única ação. \textit{EventLogger}
valida cada evento contra \texttt{uxagent\_event.schema.json}. Por fim, o
avaliador externo calcula sucesso a partir do estado final, nunca da
autodeclaração do modelo.

O prompt contém instruções operacionais gerais e a observação atual, mas não
recebe HTA, CTT, GOMS/KLM, gabarito ou estado verdadeiro. A validação do schema
é repetida localmente após a resposta do provedor.

\section{Instrumentação}
Cada linha JSONL possui versão do schema, sessão, tarefa, condição, ator,
origem sintética, versão da configuração, versão da interface, passo,
timestamp, observação, ação proposta ou executada, rejeição, suporte, erro,
oportunidade, resultado e falha técnica. Eventos incluem
\texttt{task\_start}, \texttt{observation\_emitted},
\texttt{action\_proposed}, \texttt{action\_rejected},
\texttt{action\_executed}, \texttt{support\_opened},
\texttt{task\_evaluated} e \texttt{task\_end}.

Esta instrumentação atende a rastreabilidade técnica. Entretanto, para o
entregável E3, ainda deve ser produzido um pacote versionado contendo os logs
selecionados, o dicionário de campos e os prompts e contextos efetivamente
enviados em cada etapa. O campo \texttt{obs\_id} permite correlacionar eventos.
O modo de rastreamento opcional grava prompt e contexto em arquivos separados,
sem incluir credenciais.

\section{Resultados disponíveis}
A suíte automatizada possui 70 testes aprovados. Os testes cobrem schemas,
controlador, mecanismos, sessão, interface, tarefas reservadas e execução
end-to-end. O Gemini executou T1 com sucesso em cinco propostas. Em um reteste,
T2 terminou com sucesso em sete propostas; uma repetição de
\texttt{send\_ping} foi rejeitada pelo controlador sem efeito indevido.
T3 e T4 possuem testes mock independentes. T4 foi concluída em execução real
anterior; T3 real permaneceu incompleta por repetição de ações após a remoção.

O piloto mock formal executou seis sessões, três sementes para cada tarefa de
calibração. T1 obteve \(S=1{,}00\), \(E=0\), \(A=5{,}0\) e zero uso de ajuda;
T2 obteve \(S=1{,}00\), \(E=0\), \(A=5{,}0\) e zero uso de ajuda. A diferença
de ações para a referência KLM foi \(0\) em T1 e \(-1\) em T2, cuja referência
é seis ações. Como não houve erro, \(R\) é indefinido, e não zero.

Uma amostra real do Gemini foi executada com rastreamento seguro. T1 obteve
sucesso externo com seis ações. T2 atingiu o limite de 30 passos, com quatro
ações executadas, 52 registros associados a erros e sem sucesso externo.

Esses resultados demonstram integração e robustez do controlador. Não são
resultados humanos, não estimam usabilidade e não substituem validação com
participantes.

\section{Auditoria contra o P3}
\begin{table}[t]
\caption{Estado dos requisitos principais}
\label{tab:audit}
\centering
\begin{tabularx}{\columnwidth}{p{0.34\columnwidth}X}
\toprule
Requisito & Estado\\
\midrule
Aplicação instrumentada e interface & Atendido como protótipo\\
Logs JSON estruturados e validados & Atendido\\
Mecanismos fora do prompt & Atendido\\
Agente LLM e ações restritas & Atendido\\
Verificador externo & Atendido\\
Metamensagem e paridade de signos & Parcial: documentação existe; tabela de paridade precisa explicitar cada campo observado\\
CTT com operadores do enunciado & Atendido no caderno E1 complementar\\
GOMS/KLM de especialista & Parcial: contagens existem; falta cálculo e comparação experimental\\
MoLIC por papéis com falas u:/d: & Atendido como modelo textual rastreável no caderno E1\\
Matriz MAC PP/PA/AL/RA/CE & Atendido no caderno E1 complementar\\
Ajuda sob demanda & Atendido tecnicamente; evidência de uso ainda é nula no mock\\
Piloto e dados E3 & Parcial: pacote técnico produzido; amostra real pequena\\
Relatório IEEE PDF até oito páginas & Pendente; este arquivo é a fonte \LaTeX{}\\
\bottomrule
\end{tabularx}
\end{table}

\section{Ameaças à validade}
O modelo Gemini pode repetir ações apesar do feedback; isso é registrado como
limitação observada e não é convertido em sucesso por fallback. As tarefas
T3/T4 ainda não podem calibrar prompts ou mecanismos de T1/T2. Os logs atuais
são sintéticos e a interface é um protótipo de baixa fidelidade refinado.
Também existe risco residual de paridade incompleta entre a tela visual e a
observação textual, especialmente quando novos signos forem adicionados.
Por fim, sem aprovação institucional não há recrutamento, consentimento ou
inferência sobre participantes humanos.

\section{Plano de fechamento}
Para completar o P3, devem ser executadas as seguintes ações:
\begin{enumerate}
\item converter a CTT para a notação formal exigida e produzir o diagrama
MoLIC completo, com falas do usuário e do preposto;
\item criar a matriz PP/PA/AL/RA/CE e as assinaturas MAC correspondentes;
\item ampliar o piloto real do Gemini com plano amostral pré-especificado;
\item empacotar logs, dicionário, traces e configurações como E3;
\item compilar este arquivo com IEEEtran e revisar o limite de oito páginas;
\item somente depois submeter qualquer protocolo humano à aprovação
institucional.
\end{enumerate}

\section{Conclusão}
O repositório atende a parte essencial da implementação instrumentada: há
interface funcional, acoplamento ao LLM, fronteira de informação, mecanismos
externos, logs validados e verificação independente. Portanto, a base técnica
é adequada para finalizar a entrega. Contudo, não é correto declarar que todos
os requisitos P3 foram atendidos: o PDF IEEE final ainda precisa ser fechado.
O pacote E3 e uma primeira
comparação KLM já foram produzidos para o piloto técnico, mas continuam
sintéticos e insuficientes para validação humana.

\begin{thebibliography}{10}
\bibitem{p3} Centro Federal de Educação Tecnológica, ``UXAgent-P3: Fase Desenvolvimento e Implementação do UX-Agent,'' 2026.
\bibitem{modelagem} Projeto UXAgent, ``Modelagem e artefatos de IHC,'' documentação do repositório, 2026.
\bibitem{barbosa} S. D. J. Barbosa and M. G. de Paula, ``Designing and evaluating interaction as conversation,'' DSV-IS, 2003.
\bibitem{card} S. K. Card, T. P. Moran, and A. Newell, ``The keystroke-level model for user performance time,'' \emph{Communications of the ACM}, vol. 23, no. 7, 1980.
\bibitem{dix} A. Dix, J. Finlay, G. Abowd, and R. Beale, \emph{Human-Computer Interaction}, 3rd ed. Pearson, 2004.
\bibitem{kammersgaard} J. Kammersgaard, ``Four different perspectives on human-computer interaction,'' \emph{International Journal of Man-Machine Studies}, vol. 28, 1988.
\bibitem{mackenzie} I. S. MacKenzie, \emph{Human-Computer Interaction: An Empirical Research Perspective}. Morgan Kaufmann, 2013.
\bibitem{shneiderman} B. Shneiderman and C. Plaisant, \emph{Designing the User Interface}, 5th ed. Addison-Wesley, 2010.
\end{thebibliography}

\end{document}
